\newpage
\selectlanguage{greek}
\chapter*{Περίληψη}
\addcontentsline{toc}{chapter}{Περίληψη}

Τα τελευταία χρόνια, οι ιστότοποι έχουν γίνει πιο πολύπλοκοι και ποικιλόμορφοι ως προς τον σχεδιασμό και τη δομή τους. Παρά την ποικιλομορφία αυτή, οι ιστότοποι που βρίσκονται στον ίδιο θεματικό τομέα ή εξυπηρετούν παρόμοιους σκοπούς συχνά μοιράζονται κοινά παρουσιαστικά και δομικά μοτίβα. Η αναγνώριση και η σύγκριση αυτών των μοτίβων σε διαφορετικούς τομείς ιστότοπων είναι πολύτιμη για την επίτευξη αποτελεσμάτων ως προς την ταξινόμηση ιστοσελίδων, την αυτοματοποιημένη ανάλυση διεπαφών χρήστη \textlatin{(UIs)} και την αποτύπωση των θεματικών τους τομέων. Ωστόσο, η δομή και η οπτική παρουσίαση των περιεχομένων ιστού είναι συχνά μη τυποποιημένη και ασυνεπής, αυξάνοντας τη δυσκολία τέτοιων αναλύσεων. Σε αυτήν την πτυχιακή εργασία, παρουσιάζεται το \textlatin{\textbf{S\textsuperscript{4}D}}, ένα σύστημα που αυτοματοποιεί τη διαδικασία εξαγωγής, μοντελοποίησης και σύγκρισης των παρουσιαστικών και δομικών ομοιοτήτων διαφορετικών ιστοτόπων. Το σύστημα αναλύει ένα σύνολο ιστοσελίδων μέσω τεχνικών αυτοματοποιημένης περιήγησης, μοντελοποιώντας τα \textlatin{CSS} χαρακτηριστικά των στοιχείων τους και οργανώνοντάς τα με τρόπο που επιτρέπει την περαιτέρω ανάλυση. Συγκεκριμένα, μετατρέποντας τα στοιχεία σελίδας σε διανυσματικές αναπαραστάσεις και εφαρμόζοντας τεχνικές μη-επιβλεπόμενης συσταδοποίησης, το σύστημα εντοπίζει ουσιαστικές συστάδες των στοιχείων βάσει των οπτικών τους χαρακτηριστικών. Χρησιμοποιώντας αυτήν την συσταδοποιημένη αναπαράσταση, το σύστημα μοντελοποιεί τα παραγόμενα δεδομένα για κάθε ιστότοπο, εξάγει μοτίβα διάταξης στοιχείων και υπολογίζει την ομοιότητα μεταξύ τους χρησιμοποιώντας διαφορετικές μετρικές απόστασης. Τέλος, το σύστημα παρουσιάζει έναν πίνακα ομοιότητας μεταξύ των ιστοτόπων, προσφέροντας πληροφορίες σχετικά με τη δομική και παρουσιαστική ομοιότητα των ιστοσελίδων τους.
